\documentclass[10pt]{article}

\def\Type{LessonCaelan}
\def\TypeNum{Lesson 17}
\def\TypeTitle{Antiderivatives and Indefinite Integrals }
\def\Semester{Fall 2021} %Not Used 
\def\Course{MATH 1190 \& 1210}
\def\Targets{  \bf\color{Fuchsia}I2 }

\usepackage{preambleDJ}

%%%%%%%%%%%% BEGIN DOCUMENT %%%%%%%%%%%%%%%
%%%%%%%%%%%% BEGIN DOCUMENT %%%%%%%%%%%%%%%
%%%%%%%%%%%% BEGIN DOCUMENT %%%%%%%%%%%%%%%
%%%%%%%%%%%% BEGIN DOCUMENT %%%%%%%%%%%%%%%
%%%%%%%%%%%% BEGIN DOCUMENT %%%%%%%%%%%%%%%

%%%%%%%%%%%% BEGIN DOCUMENT %%%%%%%%%%%%%%%
%%%%%%%%%%%% BEGIN DOCUMENT %%%%%%%%%%%%%%%
%%%%%%%%%%%% BEGIN DOCUMENT %%%%%%%%%%%%%%%
%%%%%%%%%%%% BEGIN DOCUMENT %%%%%%%%%%%%%%%
%%%%%%%%%%%% BEGIN DOCUMENT %%%%%%%%%%%%%%%

\begin{document}
\pagestyle{otherpages}
\thispagestyle{firstpage}
\vs{.5}

\textbf{Intended Learning Outcomes}
After this lesson, you will be able to 

\begin{itemize}
    \item Find antiderivatives of power and exponential functions. (\target{I3})
    %\item Use graph (and possibly additional information) of a function to deduce information about its antiderivative.
    \item Use algebraic simplification and/or properties of an indefinite integral to find the indefinite integral of a function. (\target{I3})
\end{itemize}

\vspace{-0.6cm}

\hrulefill\\

\begin{center}
\fbox{\it Basic Antidifferentiation Rules}
\end{center}

{\bf\color{blue} Fill in the blanks!} Let $F$ and $G$ be antiderivatives of $f$ and $g$, respectively, and let $k$ be a constant.

\begin{itemize}

\item {\bf Definition:} A function $F$ is called an {\bf antiderivative} of $f$ on an interval $I$ if \rule[-8pt]{0.5in}{0.4pt} $=$ \rule[-8pt]{0.5in}{0.4pt} for all $x$ in $I$.\\

\item A particular antiderivative of $kf(x)$ is \dotfill\rule[-8pt]{1in}{0.4pt}.\\

\item A particular antiderivative of $f(x)\pm g(x)$ is \dotfill\rule[-8pt]{2in}{0.4pt}.\\

\item The {\bf indefinite integral} of a function is the most general antiderivative of that function: 

\[\dint f(x)\,dx = \rule[-8pt]{2in}{0.4pt}\]

\end{itemize}

\

\hrulefill

\

%{\exercise} Use the definition of an antiderivative to respond to the following:\\

%Is $\dfrac{x}{x^2+1}$ an antiderivative of $\dfrac{1-x^2}{(x^2+1)^2}$? Why or why not?

%\vfill

\begin{center}
\fbox{\it Antiderivatives of Exponential Functions}
\end{center}

{\bf\color{blue}Warm-up exercise 1.} Fill in the derivative formulas: \hfill $\dfrac{d}{dx}[e^x] = \rule[-8pt]{1in}{0.4pt}$ \hfill $\dfrac{d}{dx}[a^x] = \rule[-8pt]{1in}{0.4pt}$ \hfill\,\\

{\exercise{\color{Fuchsia}I3}} Use basic antidifferentiation rules to respond to the following prompts:

\begin{itemize}
\item[(a)] Find an antiderivative of $f(x)=e^x$. Use the definition of an antiderivative to check your answer. ({\bf Hint:} you might ask yourself ``which function has a derivative of $f(x)=e^x$?")

    \vfill

\item[(b)] Find an antiderivative of $f(x)=a^x$, where $a>0$ but $a\not=1$. Use the definition of an antiderivative to check your answer. ({\bf Hint:} you might ask yourself ``which function has a derivative that contains an $a^x$ term?" to get started.)

    \vfill

\item[(c)] Based on your work above, fill in the following indefinite integral formulas:\\

\,\hfill
$\dint e^x\,dx = $\rule[-8pt]{1.5in}{0.4pt}
\hfill
$\dint a^x\,dx = $\rule[-8pt]{1.5in}{0.4pt}
\hfill\,

\end{itemize}



%%%%%%%%%
\newpage%
%%%%%%%%%

\begin{center}
\fbox{\it Antiderivatives of Power Functions \& Polynomials}
\end{center}

{\exercise{\color{Fuchsia}I3}} Find the following antiderivatives, using the space below the table to check your answers, as needed. ({\bf Hint}: you might ask yourself questions like ``which function has a derivative of $2x$?" when filling in the table.)

\renewcommand{\arraystretch}{2}

\begin{center}
    \begin{tabular}{ | c | m{2.5in} || c | m{2.5in} | }
    \hline
    Function & Particular antiderivative & Function & Particular antiderivative\\
    \hline
    $2x$ & & $6x$ & \\
    \hline
    $3x^2$ & & $x^2$ & \\
    \hline
    $\dfrac32x^{1/2}$ & & $6x^{1/2}$ & \\
    \hline
    $-x^{-2}$ & & $2x^{-2}$ & \\
    \hline
    $\dfrac1x$ & & $\dfrac{2}{x}$ & \\
    \hline
    $-\dfrac1{x^2}$ & & $\dfrac2{x^2}$ & \\
    \hline
    $x^n$ ($n\not=-1$) & & $x^{-1}$ & \\
    \hline
    \end{tabular}
\end{center}

\vfill

Now, find a particular antiderivative of $f(x)=6x-x^{-2}+6x^{1/2}+\dfrac2x-\dfrac1{x^2}$.\\

\begin{center}
$F(x) = \rule[-8pt]{5in}{0.4pt}$
\end{center}

\

Finally, find a particular antiderivative of $g(x) = x + x^{-3} + 7x^{1/3} - x^{-1} + \dfrac{1}{\sqrt[3]{x}}$.\\

\begin{center}
$G(x) = \rule[-8pt]{5in}{0.4pt}$
\end{center}

\

Based on your work above, fill in the following indefinite integral formulas:\\

\,\hfill
$\dint x^n\,dx = \underset{\text{assuming }n\not=-1}{\rule[-8pt]{1.5in}{0.4pt}}$
\hfill
$\dint \frac1x\,dx = $\rule[-8pt]{1.5in}{0.4pt}
\hfill\,

\

%%%%%%%%%
\newpage%
%%%%%%%%%

\begin{center}
\fbox{\it Summary of Indefinite Integrals}
\end{center}

{\bf\color{blue} Fill in the blanks!} Assume $k$ is any constant. Furthermore, assume $f$ has antiderivative $F$ and $g$ has antiderivative $G$. Fill in the blanks below to complete the table containing some basic indefinite integrals.

\begin{center}
{\it Integral of a Constant}\\[0.1in]
\begin{tabular}{ r c l}
$\displaystyle\int k\, dx$ & $=$ & \rule[-8pt]{1.75in}{0.4pt}
\end{tabular}
\end{center}

\

\begin{center}
{\it Integral of Power Functions}\\[0.1in]
\begin{tabular}{ r c l r c l}
$\displaystyle\int x^n\, dx$ & ${=}$ & $\underset{\text{assuming }n\not=-1}{\rule[-8pt]{1.75in}{0.4pt}}$ & $\displaystyle\int\frac1x\,dx$ & $=$ & \rule[-8pt]{1.75in}{0.4pt}\\[0.1in]
\end{tabular}
\end{center}

\

\begin{center}
{\it Integral of Exponential Functions}\\[0.1in]
\begin{tabular}{ r c l r c l}
$\displaystyle\int e^x\, dx$ & $=$ & \rule[-8pt]{1.75in}{0.4pt} & $\displaystyle\int a^x\,dx$ & $=$ & $\underset{\text{assuming }a>0\,\&\,a\not=1}{\rule[-8pt]{1.75in}{0.4pt}}$\\[0.1in]
\end{tabular}
\end{center}

\

\begin{center}
{\it Linearity Properties of Indefinite Integrals}\\[0.1in]
\begin{tabular}{ r c l r c l}
$\displaystyle\int \big( f(x) \pm g(x) \big)\, dx$ & $=$ & \rule[-8pt]{1.75in}{0.4pt} & $\displaystyle\int kf(x) \,dx$ & $=$ & \rule[-8pt]{1.75in}{0.4pt}\\[0.1in]
\end{tabular}
\end{center}

\

{\exercise{\color{Fuchsia}I3}} Find the following indefinite integrals. ({\bf Hint:} when you can't find an answer using the rules above, try rewriting and/or simplifying algebraically first!) 

\begin{enumerate}

\item[(a)] $\dint \sqrt2 \,dx$

\vfill

\item[(b)] $\dint \left(\sqrt[3]{u}+4u^{-1}+\frac{1}{u^2}\right)\,du$

\vfill

\item[(c)] $\dint \left(\dfrac{y^4-1}{y^2} \right)\,dy$ 

\vfill
\vfill

\end{enumerate}

%%%%%%%%%
\newpage%
%%%%%%%%%

{\bf\color{blue}Extra Practice 1.}\ltrm{\color{Fuchsia}I3} Find the following indefinite integrals:

\begin{enumerate}

\item[(a)] $\dint \left(2 - 6t^3 + \dfrac{3}{t^2} - e^t \right)\,dt$

\vfill

\item[(b)] $\dint \dfrac{x^3+x^2-x+1}{2x^2}\,dx$

\vfill
\vfill

\item[(c)] $\dint \left( z^2 + \dfrac1{z^2}\right) \left( z - \dfrac{1}{z} \right) \, dz$

\vfill
\vfill

\end{enumerate}

\end{document}


%%%%%%%%%%%%%%%%%%%%%%%%%
Good! Now, use what you know about antiderivatives and derivative rules to fill in the following blanks, again marking a blank with an ``$X$" if the blank cannot be filled in using the information we have so far.

\begin{itemize}

\item A particular antiderivative of $\dfrac{f(x)}{g(x)}$ is \dotfill\rule[-8pt]{2in}{0.4pt}.\\

\item A particular antiderivative of $\dfrac{lo(x)hi'(x) - hi(x)lo'(x)}{[lo(x)]^2}$ is \dotfill\rule[-8pt]{2in}{0.4pt}.\\

\item A particular antiderivative of $f(x)g(x)$ is \dotfill\rule[-8pt]{2in}{0.4pt}.\\

\item A particular antiderivative of $f'(x)g(x)+f(x)g'(x)$ is \dotfill\rule[-8pt]{2in}{0.4pt}.

\item A particular antiderivative of $f(g(x))$ is \dotfill\rule[-8pt]{2in}{0.4pt}.\\

\item A particular antiderivative of $f'(g(x))\cdot g'(x)$ is \dotfill\rule[-8pt]{2in}{0.4pt}.\\

\end{itemize}